\section{Two all-iteration counterexamples}

\begin{proposition}[Small two-cycle]\label{prop:small}
Let $K(u)=e^{-u^2/4}(u^4+22u^2+241)$ and
$g(t)=K(\sqrt t)$. Put $A(t)=t^2+22t+241$
and $B(t)=t^2+22t+1$. Then, for all $n\ge0$
and $t\ge0$,
\[
\mathcal T^ng(t)=2^{2^n-1}e^{-2^nt/4}
\begin{cases}A(t),&n\text{ even},\\B(t),&n\text{ odd}\end{cases}>0.
\]
Moreover $(\log K)''<-1/4$ for every real $u$, yet
\[
\frac{1}{2\sqrt\pi}\int_{\mathbb R}K(u)e^{izu}du
=e^{-z^2}\big((4z^2-17)^2+8\big)
\]
has four nonreal zeros.
\end{proposition}
\begin{proof}
Direct differentiation shows
$\mathcal TA=2B$ and $\mathcal TB=2A$. Since
$\mathcal T(ce^{-at}P)=c^2e^{-2at}\mathcal TP$,
the iterate formula follows inductively.
With $Q(u)=u^4+22u^2+241$, one has
\begin{align*}
\frac{Q^2}4-QQ''+(Q')^2
&=\frac{u^8}4+15u^6+\frac{571u^4}2\\
&\quad+727u^2+\frac{15665}4>0.
\end{align*}
Thus $(\log K)''=-1/4-
\bigl(Q^2/4-QQ''+(Q')^2\bigr)/Q^2<-1/4$.
Gaussian Fourier differentiation gives the claimed transform;
its zeros satisfy $z^2=(17\pm2\sqrt2 i)/4$.
\end{proof}

\begin{proposition}[Original rational-root construction]
Let
\[
K_*(u)=e^{-u^2/4}(62161+1848u^2+16u^4).
\]
Set $A_*(t)=62161+1848t+16t^2$ and
$B_*(t)=44561+1848t+16t^2$.
For $g_*(t)=K_*(\sqrt t)$,
\[
\mathcal T^n g_*(t)=32^{2^n-1}e^{-2^nt/4}
\begin{cases}A_*(t),&n\text{ even},\\
B_*(t),&n\text{ odd}\end{cases}>0.
\]
The kernel satisfies $(\log K_*)''<-3/8$ on $\mathbb R$,
but its normalized Fourier transform is
\[
F_*(z)=e^{-z^2}\big[(16z^2-255)^2+1024\big],
\]
whose zeros are exactly $z=\pm4\pm i/4$.
\end{proposition}
\begin{proof}
One checks $\mathcal TA_*=32B_*$ and
$\mathcal TB_*=32A_*$. The exponential scaling identity
used above proves the iteration formula. With
$Q_*(u)=62161+1848u^2+16u^4$,
\[
Q_*(u)/8-Q_*''(u)=32593/8+39u^2+2u^4>0
\]
implies the stated log-concavity bound.
Fourier differentiation gives $F_*$; since
$16(4+i/4)^2-255=32i$, it has the claimed root,
and evenness and complex conjugation supply the other three.
\end{proof}

The idea of counterexamples based on positive
Gaussian-polynomial kernels predates this note; see
Csordas~\cite{Cs}, Example 3.12. No priority claim is made
for that general construction.
